%% Constraint_Is_Not_a_Limitation_APA7.tex
%% APA 7.0 manuscript format
%% Géoui & Huzard (2026)
%%
%% Compile sequence:
%%   pdflatex Constraint_Is_Not_a_Limitation_APA7
%%   biber    Constraint_Is_Not_a_Limitation_APA7
%%   pdflatex Constraint_Is_Not_a_Limitation_APA7
%%   pdflatex Constraint_Is_Not_a_Limitation_APA7
%% Or simply: latexmk -pdf Constraint_Is_Not_a_Limitation_APA7

\documentclass[man,12pt,donotrepeattitle]{apa7}

\usepackage[utf8]{inputenc}
\usepackage[T1]{fontenc}
\usepackage[american]{babel}
\usepackage{csquotes}
\usepackage[style=apa,sortcites=true,sorting=nyt,backend=biber]{biblatex}
\DeclareLanguageMapping{american}{american-apa}
\addbibresource{references.bib}
\usepackage{hyperref}
\hypersetup{colorlinks=true,linkcolor=blue,citecolor=blue,urlcolor=blue}

%% ── Document metadata ────────────────────────────────────────────────────

\title{Constraint Is Not a Limitation: Graphs, Ontologies, and the Future of Efficient LLM Systems in Life Sciences}

\author{Thibault G\'{e}oui \and Damien Huzard}
\affiliation{Independent Scholars}
\leftheader{CONSTRAINT IS NOT A LIMITATION}

\authornote{%
Thibault G\'{e}oui, Independent Scholar.
Damien Huzard, Metadatapp.
Correspondence concerning this article should be addressed to Damien Huzard.
Email: \href{mailto:damien@metadatapp.net}{damien@metadatapp.net}
}

\keywords{knowledge graphs, ontologies, large language models,
retrieval-augmented generation, life sciences, FAIR data,
knowledge infrastructure, hallucination reduction, hybrid architecture}

%% ── Abstract ─────────────────────────────────────────────────────────────

\abstract{%
The rapid adoption of large language models has created the impression
that structural knowledge infrastructure---taxonomies, ontologies,
metadata standards, and knowledge graphs---may no longer be necessary.
If a model can read anything and reason across it, why invest in the
painstaking work of building and maintaining formal knowledge systems?

This paper argues that impression is economically and technically
misleading, and that its cost is becoming visible. LLM inference is not
free infrastructure: it carries variable compute, energy, and governance
costs that compound at scale, and the pricing models that made early
adoption feel frictionless are already shifting toward usage-sensitive
billing. At the same time, the technical limitations of unstructured LLM
use---context degradation, retrieval fragmentation, hallucination risk,
and the absence of provenance---do not disappear as models grow more
capable. Many of these limitations are not solely capability problems;
they are also architectural, epistemic, and governance problems---structural
challenges that persist regardless of model size because they arise from
the absence of shared knowledge infrastructure, not from any failure of
model reasoning. The cost of this absence takes five recurring forms:
conceptual waste, retrieval waste, validation waste, human waste, and
energy waste, each compounding at scale in ways that structured knowledge
infrastructure can systematically address.

The solution is not to abandon LLMs. It is to position them correctly
within a hybrid architecture in which metadata standards, ontologies,
knowledge graphs, deterministic validators, and human expertise each
perform the role they are best suited for, and the model is invoked for
synthesis and ambiguity management rather than for tasks that structured
systems handle more reliably and cheaply.

We argue that this architectural shift also has an organizational
dimension that has received insufficient attention. The people needed to
build and sustain knowledge infrastructure---domain experts who understand
biology, chemistry, and therapeutic context and can work with ontologies,
metadata schemas, and validation rules---are not the same people who build
AI models. Organizations in the life sciences that do not plan for this
talent will find themselves increasingly dependent on AI vendors not only
for compute, but for interpretation.

The paper is positioned as a bridge between two audiences that rarely
read the same documents: the technical architects, knowledge engineers,
and informatics leads who build and operate AI-supported data systems,
and the senior decision-makers---research leaders, data science directors,
digital transformation teams, heads of R\&D informatics---who set strategy
and approve investment. The intent is to make the technical case legible
to leadership and the strategic case actionable for architects, so that
the same evidence supports both budget decisions and design decisions. Its
central claim is that constraint is not a limitation. It is how
organizations stay in control of cost, quality, and knowledge.
}

%% ─────────────────────────────────────────────────────────────────────────
\begin{document}
\maketitle

%% ── Section 1 ────────────────────────────────────────────────────────────

\section{Section 1 --- The Illusion of ``Just Use AI''}

\subsection{The Moment Things Changed}

For anyone who had spent years trying to convince a leadership team that
investing in taxonomies, ontologies, and data foundations was worth the
effort, the arrival of large language models felt like a gut punch.

Suddenly, you could take an article, a clinical report, a protocol
document, and run it through a model. Ask it to extract concepts, map
relationships, summarize findings. And it worked remarkably well, and
without any of the painstaking upfront work that structured knowledge
management had always required. No taxonomy team. No controlled
vocabulary. No months-long ontology engineering project. Just a prompt,
and something that looked a lot like intelligence.

For a while, a serious question hung in the air: \textit{do we actually
need taxonomies and ontologies anymore?}

It was not a naive question. It was the right question to ask. It did not
have a clean answer, until the economics started to shift.

\subsection{The Subsidized Era}

What the early LLM experience produced was not frictionless intelligence.
It produced the \textit{impression} of frictionless intelligence, at a
price that did not reflect reality.

The AI industry in its early scaling phase operated, whether intentionally
or not, with a logic familiar from other platform markets: make the product
so cheap and so compelling that it becomes embedded in workflows before
anyone has fully priced the dependency. One industry analyst described the
dynamic bluntly: AI products have significant variable costs because every
API call and token processed adds cost, and those costs were not being
passed on \parencite{a16z2024}. The infrastructure underpinning each model
call---chips, energy, data centers, inference capacity---is closer to
industrial infrastructure than to classic software \parencite{reutersbreaking2025}.
It does not scale for free.

This does not mean prices will move in only one direction. Inference costs
per token have fallen dramatically and continue to do so, but the rate
varies enormously across performance tiers: between 9$\times$ and
900$\times$ per year depending on the task, a hundred-fold difference in
the pace of cost reduction \parencite{epochai2025}. The more important
observation is about \textit{total enterprise exposure}. Unit costs may
fall while overall spend rises, as reasoning-capable models are deployed
deeper into workflows, context windows grow, and the number of embedded AI
processes multiplies. An organization that has restructured its data
pipelines, its search, its literature monitoring, and its regulatory
workflows around a single model provider has not reduced its dependency by
watching the per-token price fall. It has increased it.

The transition from experimentation to exposure is already visible in how
AI vendors are repricing their products. Early subscription models made AI
feel like a flat-cost utility: bounded, predictable, absorbing whatever
usage organizations threw at it. Agentic workflows break that assumption
entirely. A human user generates dozens of prompts per day; an autonomous
workflow can generate thousands of model calls, intermediate reasoning
steps, retries, and validations in the same period. The underlying cost
structure of AI---inference, context processing, output generation, and
repeated interpretation---is becoming visible to the organizations that
consume it. GitHub Copilot, for example, is migrating from request-based
pricing to AI Credits billed by token consumption, making the cost of each
operation explicit \parencite{rodriguez2026}. Anthropic has separately
begun putting outside agent tool usage behind a separate credit meter,
distinct from standard subscription access, a direct recognition that
automated workflows consume compute at a fundamentally different rate than
human users \parencite{fried2026}. Reports indicate that some enterprise
customers, including ServiceNow and Uber, burned through their entire
annual AI token budgets before the year was out \parencite{fried2026}.
These are not isolated product decisions. They are the moment the
underlying economics become legible.

\subsection{What the Model Cannot Do on Its Own}

Behind the apparent ease of the LLM interaction lies a technical reality
that matters for anyone designing systems at scale.

When a model receives a long document, or ten documents, or a hundred, it
does not reliably use all of the information it has been given. Research
has shown that model performance can degrade when relevant information is
located in the middle of a long context, a phenomenon sometimes called
``lost in the middle'' \parencite{liu2024}. Extending the context window
does not fully solve this: only a limited number of state-of-the-art models
maintain consistent accuracy beyond 64,000 tokens, and the computational
cost of doing so rises steeply \parencite{leng2024}.

The model must also, in the absence of structure, reconstruct meaning from
scratch each time. It has to infer which terms are equivalent, which
measurement units are compatible, which source is authoritative, which
metadata is missing. None of that work is free. It is paid for in tokens,
in compute, and in error rates. Retrieval-augmented generation improves
accuracy for knowledge-intensive tasks by combining the model with external
knowledge sources \parencite{gao2024}, but standard retrieval returns text
fragments without preserving the relationships between them, leaving the
model to infer connections that could have been made explicit.

The result is a set of costs that compound at scale: \textit{conceptual
waste}, when the model re-infers relationships that were already known;
\textit{retrieval waste}, when irrelevant or redundant content fills the
context; \textit{validation waste}, when outputs must be corrected because
the input lacked structure; \textit{human waste}, when experts spend time
reviewing generated text that a deterministic system could have checked;
and \textit{energy waste}, because inference energy correlates directly
with output token length and model size \parencite{poddar2025,dauner2025}.
As a concrete instance we develop in Section 8, consider the construction
of a virtual control group from historical preclinical data: a task in
which the cost of asking a model to infer comparability across species,
strain, route of administration, and endpoint definition from unstructured
text---without standardized metadata or ontology-mapped variables---becomes
immediately visible, both in inference cost and in the regulatory
acceptability of the result.

\subsection{A Question of Architecture}

The relevant metric is not price per million tokens. It is tokens per
useful decision; more broadly, the total organizational cost of producing
a reliable, auditable, reusable output.

That reframing changes the question. It is not: \textit{should we use AI?}
The answer to that is clearly yes. It is: \textit{which AI system, for
which task, with which supporting infrastructure, at what cost, and with
what governance?}

Retrieval-augmented generation with a cost-optimized routing strategy can
maintain comparable performance to long-context inference at a fraction of
the compute \parencite{li2024}. Deterministic validation systems can catch
errors that would otherwise require model calls. Smaller models, guided by
structured knowledge, can handle tasks that currently route unnecessarily
to frontier systems.

This is a question of architecture. And architecture, in turn, is a
question of what knowledge infrastructure exists before the model is ever
invoked.

The organizations that will manage AI costs effectively, and that will be
least exposed when pricing normalizes, are not necessarily the ones with
access to the most powerful models. They are the ones that have invested in
making their knowledge machine-actionable before it reaches the model:
through metadata standards, ontologies, knowledge graphs, and validation
layers. Infrastructure that reduces the model's interpretive burden.
Infrastructure that the model cannot provide for itself.

\subsection{Scope and Decision Horizon}

A note on the scope of the argument is warranted before we proceed. The
case we make in this paper applies to a 3--5 year decision horizon: the
window over which most current life science and pharmaceutical organizations
are committing to AI infrastructure, hiring plans, vendor contracts, and
regulatory positioning. We are not making an indefinite claim about what
large language models will or will not be able to do at some future date.
We are making a claim about which architectural and organizational choices
are defensible \textit{given current evidence} about model behaviour,
\textit{current and announced} pricing trajectories, and \textit{current}
regulatory expectations around auditability and provenance in regulated
science. Section 9 addresses explicitly the conditions under which the
argument would weaken, and what evidence would tell us we are approaching
them.

The rest of this paper explains what that infrastructure is, why it matters
technically, what it takes to build and sustain it organizationally, and
why the investment case is stronger now than it has ever been.

%% ── Section 2 ────────────────────────────────────────────────────────────

\section{Section 2 --- What LLMs Cannot Do Efficiently Alone}

\subsection{Intelligence Without Memory}

A large language model is, at its core, a pattern-completion and reasoning
system of remarkable capability. It can read a dense scientific paper and
produce a lucid summary. It can identify relationships between entities
across a long document. It can translate, reformat, compare, and explain.
These are genuine capabilities, and they have changed what is possible in
knowledge work.

What a large language model is not, however, is a knowledge management
system. It has no persistent memory of your organization's data. It does
not automatically know which terms in your corpus are equivalent, which
measurement units are compatible across studies, which source should be
treated as authoritative, or which metadata fields are missing from a
record. Without structure provided from outside, it must reconstruct all
of that context from text, every time it is invoked.

This distinction matters because it is the source of a class of costs
that do not disappear as models become more capable. These are not solely
capability costs; they are also architectural, epistemic, and governance
costs---they arise from the absence of structured knowledge made available
to the model, not from any failure of model reasoning. A more powerful
model reconstructing context it was never given is still reconstructing
context it was never given.

\subsection{The Five Forms of Waste (Proposed Framework)}

When unstructured text is the primary input to an LLM-based workflow, we
observe that several categories of waste compound at scale. In this section
we propose a working taxonomy that groups them into five categories. We
present it as an analytic framework, not as an established or empirically
validated typology: the categories are organizationally useful for thinking
about where structured knowledge infrastructure earns its return, but the
boundaries between them are not perfectly mutually exclusive (e.g.,
conceptual confusion downstream produces validation waste; retrieval
failures often cascade into human waste), and we are not claiming that the
five categories are collectively exhaustive of all forms of inefficiency in
unstructured LLM workflows. The magnitudes attached to each category in any
given deployment are also organization-specific and depend on task mix,
model choice, and workflow design; we present qualitative descriptions and
cite representative evidence where available, rather than claiming a fixed
quantitative ranking.

With those caveats, we find the five categories below useful in practice
for diagnosing where structured infrastructure pays back.

\textit{Conceptual waste} occurs when a model re-infers relationships that
are already known and stable. If your organization has established that
``locomotor activity,'' ``distance travelled,'' ``velocity,'' and
``home-cage monitoring'' are related but not interchangeable concepts in
preclinical research, that distinction should not need to be reconstructed
from raw text on every query. An ontology encodes it once. Without one,
the model guesses, sometimes correctly, sometimes not, and always at a cost.

\textit{Retrieval waste} occurs when irrelevant or redundant content fills
the context window. Standard retrieval-augmented generation returns text
chunks ranked by semantic similarity. It does not preserve the relationships
between those chunks, the provenance of the information they contain, or the
conditions under which a finding was produced. The model receives fragments
and must synthesize them without knowing how they connect. Knowledge
graph-guided retrieval addresses this directly: instead of returning
isolated passages, it retrieves entities, relations, provenance, and
neighboring evidence as a compact, structured bundle
\parencite{zhu2025,peng2025}.

\textit{Validation waste} occurs when outputs must be corrected after
generation because the inputs lacked structure. A model asked to extract
experimental parameters from a document without a metadata schema has no
reference point against which to check its own output. Errors surface
downstream, often after human review, occasionally after a decision has
already been made. A validation layer defined upstream, using standards
such as SHACL for RDF graphs or JSON Schema for structured records, catches
structural errors before they reach the model and flags generated outputs
that violate known constraints \parencite{wilkinson2016}.

\textit{Human waste} occurs when subject matter experts are positioned at
the end of the pipeline to repair outputs that a deterministic system could
have caught. This is the most expensive form of waste, not only in salary
terms, but in the opportunity cost of expert attention. A scientist
reviewing a list of incorrectly extracted entity relationships is not doing
science. The correct position for expert judgment is upstream: defining the
conceptual system, validating the schema, resolving genuinely ambiguous
cases, and reviewing high-stakes outputs. Not correcting routine formatting
errors or obvious hallucinations that a constraint layer would have
prevented.

\textit{Energy waste} occurs because every token processed has a
computational and environmental cost. Inference energy correlates directly
with output token length and response time \parencite{poddar2025}. Model
size and reasoning depth drive emissions further \parencite{dauner2025}.
General-purpose generative AI systems can be orders of magnitude more
energy-intensive than task-specific systems for the same underlying task
\parencite{luccioni2024}. Sending forty pages to a frontier model and asking
it to find the relevant paragraph is not only expensive in monetary terms;
it is a poor engineering choice when a graph query could retrieve the
relevant subgraph in milliseconds.

\subsection{The Long Context Problem}

One apparent solution to the context governance challenge is simply to
extend the context window. If the model can read everything, the problem of
deciding what to include disappears.

This solution is technically weaker than it appears. Research has
demonstrated that model performance can degrade depending on where relevant
information appears in a long prompt; models do not reliably attend to
content located in the middle of a long context, a finding that holds even
for models specifically designed for long-context tasks \parencite{liu2024}.
Extending the context window to hundreds of thousands of tokens does not
fully resolve this: only a limited number of state-of-the-art models
maintain consistent accuracy beyond 64,000 tokens, and performance variation
across models remains substantial \parencite{leng2024}.

The computational cost of long-context inference also rises steeply.
Standard transformer attention scales quadratically with sequence length,
and although a substantial body of work---efficient attention variants such
as FlashAttention, linear and sparse attention, sliding-window
architectures, and mixture-of-experts schemes---has reduced the effective
constant factor and in some cases the asymptotic complexity, long-context
inference remains a significant engineering and cost challenge in production
at scale \parencite{tay2022,dao2022}. The point of this paper does not
depend on the precise complexity class: even with the most efficient current
implementations, retrieval-augmented approaches with structured retrieval
maintain comparable task performance at a fraction of the compute, with
hybrid routing between RAG and long-context inference offering a practical
middle path \parencite{li2024}.

The relevant metric, therefore, is not how large a context window a model
supports. It is how precisely the right information can be identified,
structured, and delivered to the model before the context is assembled. That
is primarily a knowledge infrastructure problem---not solely a model
capability problem.

\subsection{What the Model Needs, but Cannot Provide Itself}

Stated directly: an LLM is good at language. It is not good at being a
database, a schema validator, a provenance tracker, or an ontology. Asking
one system to do all four jobs simultaneously produces a system that does
all four jobs expensively and inconsistently.

The research literature on this point is no longer speculative. Knowledge
graphs are a promising external-knowledge source for reducing hallucinations
and improving reasoning accuracy \parencite{agrawal2024}, though the
integration of KG and LLM systems remains an active area with unresolved
benchmark and evaluation challenges \parencite{lavrinovics2024}.
Ontology-grounded retrieval, in which a domain ontology is used to retrieve
a minimal and conceptually coherent context rather than a large
undifferentiated one, has demonstrated improvements in fact recall,
attribution speed, and fact-based reasoning relative to standard retrieval
baselines \parencite{sharma2024}.

The practical implication is not that LLMs should be avoided. It is that
they should be positioned correctly in a larger system: invoked for
synthesis, explanation, and the management of genuine ambiguity, after
structured retrieval, ontology mapping, and validation have done what they
are better suited to do.

The next section describes what that infrastructure consists of and how
each component contributes.

%% ── Section 3 ────────────────────────────────────────────────────────────

\section{Section 3 --- What Graphs, Ontologies, and Metadata Actually Do}

\subsection{Four Components, Four Jobs}

There is a persistent tendency to treat ``structured knowledge'' as a
single monolithic concept, as if metadata, ontologies, and knowledge graphs
were interchangeable terms for the same thing. They are not. Each component
solves a different problem, and conflating them leads to poor architectural
decisions.

A \textit{metadata standard} defines what must be described about a piece
of data. Who produced it, under which conditions, with which variables,
units, protocol, species, model, instrument, version, and license. Without
this, a dataset is interpretable only by the person who created it, and
only for as long as they remember its context. FAIR data principles, first
articulated in 2016 and now foundational to scientific data governance, make
explicit that machine-actionability is the goal: data and metadata should
be findable, accessible, interoperable, and reusable by machines, not only
by humans \parencite{wilkinson2016}. In clinical and translational research
the same logic has produced mature, widely adopted instances: the
Observational Medical Outcomes Partnership Common Data Model (OMOP CDM),
maintained by the OHDSI community, standardizes the representation of
observational health data so that analyses can be executed identically
across sites and institutions \parencite{hripcsak2015}, and HL7 FHIR
provides a standards-based framework for the structured exchange of clinical
information across systems \parencite{bender2013}. These are not
aspirational: both are operational infrastructure in regulated clinical
research today, and they exemplify what \textit{machine-actionable} metadata
at the data-layer level looks like in practice. Metadata is not
administrative paperwork. It is the control surface through which automated
systems can find, filter, validate, and combine data without human mediation.

An \textit{ontology} defines what things mean and how they relate. It is a
formal, shared specification of concepts, their boundaries, their synonyms,
and the relationships between them. In preclinical research, for example,
``locomotor activity'' may be related to ``movement,'' ``distance
travelled,'' ``velocity,'' ``immobility,'' and ``home-cage monitoring,''
but these are not interchangeable terms. The differences matter for data
integration, for comparison across studies, and for any automated system
that needs to reason about them. An ontology formalizes those distinctions
in a way that both machines and domain experts can use. Seen from an AI
efficiency perspective, an ontology is a form of semantic compression: a
stable identifier and a defined relation can replace paragraphs of repeated
textual explanation, reducing the interpretive burden on any model that
consumes the data \parencite{sharma2024,batista2022}.

A \textit{knowledge graph} stores entities and their relationships
explicitly, rather than leaving them buried in natural language text.
Instead of a paragraph saying ``mouse A received treatment B after surgery C
on date D,'' the graph represents: Animal $\rightarrow$ received
$\rightarrow$ Treatment; Treatment $\rightarrow$ occurredAfter $\rightarrow$
Surgery; Surgery $\rightarrow$ hasDate $\rightarrow$ D; Measurement
$\rightarrow$ belongsTo $\rightarrow$ Animal. This representation is
directly queryable. It supports targeted retrieval of specific entity
relationships, provenance chains, neighboring evidence, and structured
comparisons across studies, without requiring a model to read and parse
unstructured text to reconstruct those connections \parencite{callahan2024,ma2025}.

A \textit{validation layer} checks whether data and generated outputs
conform to the expected structure. Standards such as SHACL provide a formal
language for validating RDF graphs against defined conditions. OWL enables
rich representation of knowledge about things and their relationships in a
computationally usable form. JSON-LD provides a JSON-compatible way to
serialize linked data, making semantic structure easier to integrate into
existing web and API systems. Complementary packaging conventions such as
RO-Crate make it possible to bundle data, metadata, code, and provenance
into a portable, machine-actionable research object that can be validated
and audited as a unit \parencite{soilandreyes2022}. Validation is what makes
the entire pipeline auditable: it creates a legible record of what was
checked, what passed, and what was flagged, independent of the model's
confidence.

\subsection{What Each Component Contributes to LLM Efficiency}

The value of this infrastructure is not abstract. Each component reduces
a specific category of the waste described in Section 2.

Metadata standards reduce \textit{retrieval waste} and \textit{validation
waste} by making datasets filterable before they reach the model, and by
providing a schema against which outputs can be checked. Machine-actionable
metadata templates that encode domain-specific requirements have been shown
to support verifiable metadata quality, modularity, and interoperability
across systems \parencite{batista2022,musen2022}.

Ontologies reduce \textit{conceptual waste} by encoding stable distinctions
once, rather than requiring models to reconstruct them on each query. They
also guide extraction: ontology-based recommendations have been shown to
help users enter metadata more rapidly and accurately, suggesting that
constraints reduce human burden rather than increase it
\parencite{martinezromero2017}. Portable, embeddable metadata authoring
tools that bring template-driven, ontology-guided entry into the systems
where researchers already work---rather than requiring them to switch into
a dedicated curation environment---reduce the friction of producing
standards-aligned metadata at the point of capture \parencite{oconnor2026}.

Knowledge graphs reduce \textit{retrieval waste} by enabling graph-based
retrieval that preserves entity relationships and provenance rather than
returning semantically similar but structurally disconnected fragments.
KG-guided RAG uses graph relationships to expand and organize retrieved
content, improving both retrieval and response quality \parencite{zhu2025}.
Biomedical knowledge graph ecosystems have demonstrated this at scale,
integrating ontologies and heterogeneous life science data to support
AI-powered research across large corpora \parencite{callahan2024}.

Validation layers reduce \textit{human waste} and \textit{validation waste}
by catching errors deterministically, before they require expert review. A
system that validates outputs against a known schema surfaces problems at
the point of generation rather than at the point of downstream use, where
the cost of correction is highest.

\subsection{The Result: A Compact Evidence Bundle}

The practical outcome of combining these components is a workflow in which
the model receives not a large undifferentiated context, but a compact,
structured evidence bundle: the relevant entities, their defined
relationships, the sources from which they were drawn, the conditions under
which the findings were produced, and a clear flag for anything that remains
genuinely ambiguous or unresolved.

This is what the model is actually good at working with. Not forty pages of
raw text. A precisely assembled question, with precisely assembled evidence,
and a clearly scoped task.

The naive workflow asks: ``Here are forty pages. Read everything. Understand
the domain. Find what is relevant. Infer the structure. Answer my
question.''

The structured workflow asks: ``Here is the question. Query the graph.
Retrieve the relevant entities, relations, evidence, and provenance. Send
the compact subgraph and any unresolved ambiguity to the model.''

The difference in token consumption, latency, error rate, and auditability
is not marginal.

%% ── Section 4 ────────────────────────────────────────────────────────────

\section{Section 4 --- Why Structure Reduces Cost, Hallucination, and Energy}

\subsection{The Hallucination Problem Is Partly an Architecture Problem}

Hallucination in LLM outputs is often framed as a model quality problem,
something to be solved by better training, more parameters, or improved
prompting. That framing is incomplete. A substantial portion of
hallucination risk is architectural: it arises when a model is asked to
synthesize claims without being given the structured evidence needed to
ground them.

Standard RAG improves on pure generative output by providing retrieved text
as context, but it leaves a significant gap. The model still receives
fragments without relationship information, without explicit provenance, and
without a validation layer that could flag inconsistencies. It can mix
incompatible facts from different conditions, overgeneralize from limited
evidence, or ignore the source metadata that would make a finding
contextually meaningful \parencite{lavrinovics2024}.

Graph-based retrieval imposes stronger constraints. A well-constructed
GraphRAG workflow can specify which entities are valid, which relationships
hold, which source supports each claim, which experimental conditions apply,
which claims are directly observed versus inferred, and which answer
components lack supporting evidence \parencite{peng2025,barry2025}.
Hallucination does not become impossible, but it becomes easier to prevent,
detect, and audit, because the system can compare generated claims against
explicit structured evidence rather than relying solely on the model's
internal confidence.

The research survey evidence is appropriately cautious on this point.
Knowledge graphs are a promising external-knowledge source for reducing
hallucinations and improving reasoning accuracy, but results are
task-specific and benchmarks for KG-LLM integration remain inconsistent
\parencite{agrawal2024,lavrinovics2024}. The claim is not that graphs
eliminate hallucination. It is that they give you the tools to catch it.

\subsection{Token Reduction Is Energy Reduction}

There is a dimension of this argument that extends beyond organizational
cost and into environmental responsibility, particularly relevant for large
pharmaceutical and life science organizations with public sustainability
commitments.

Inference energy correlates directly with output token length and response
time \parencite{poddar2025}. Reasoning-capable models and larger models
incur substantially higher emissions per query \parencite{dauner2025}.
General-purpose generative AI is orders of magnitude more energy-intensive
than task-specific systems for many tasks \parencite{luccioni2024}, and
carbon accounting for LLMs requires lifecycle thinking that includes
training, hardware, operational energy, and inference through user-facing
APIs \parencite{luccioni2023}.

The practical implication is that token economy and energy economy are the
same problem. A structured workflow that reduces context size, enables
caching of stable concepts and schemas, and routes work away from frontier
models where smaller models or deterministic systems suffice, is not only
cheaper to operate. It is more defensible from a sustainability standpoint.
Energy per token should complement accuracy benchmarks as a model selection
criterion, and model selection and reasoning depth can be routed dynamically
to balance accuracy and energy \parencite{wilhelm2025}.

Data centre electricity demand is growing strongly, with AI-focused
infrastructure growing particularly fast \parencite{noffsinger2025}. At
system scale, unnecessary model calls are not neutral. They are a governance
question as much as an engineering one.

\subsection{The Economic Argument for Routing}

The architectural logic of the paper is ultimately an economic one. Not
every task requires the same level of intelligence, and pricing AI workflows
as if they do is a form of systematic waste.

Do you need a frontier model to rewrite an email? To check whether a
metadata field matches a controlled vocabulary? To verify that an extracted
entity name matches a known identifier in an ontology? No. These are
deterministic or near-deterministic tasks. Routing them to a frontier model
because it is convenient is the equivalent of using a surgical robot to open
a package.

Deterministic systems should handle what is deterministic. Graphs should
handle relationships. Ontologies should handle meaning. LLMs should handle
synthesis, explanation, and ambiguity. Humans should handle expertise and
accountability. The right architecture uses each component at the point
where its cost-to-value ratio is best \parencite{ma2025,wilhelm2025}.

This is not a theoretical position. It is increasingly the practical
experience of organizations that have deployed AI at scale and discovered
that the initial convenience of routing everything to a single large model
does not survive contact with real operational costs.

%% ── Section 5 ────────────────────────────────────────────────────────────

\section{Section 5 --- A Practical Architecture for Efficient LLM Workflows}

\subsection{The Pipeline}

A concrete architecture for efficient LLM-based workflows in
knowledge-intensive domains follows a sequence that separates deterministic
processing from model-dependent processing, and positions validation at both
ends.

Raw documents, datasets, protocols, and papers enter a deterministic parsing
stage. Metadata is extracted and mapped against known schemas. Entities are
identified and mapped to ontology terms. The resulting structured output is
validated against formal constraints, using SHACL, JSON Schema, or
equivalent standards depending on the representation format. Validated
records are stored in a knowledge graph with full provenance.

When a query arrives, it is answered through graph query or hybrid
retrieval: the system retrieves the relevant entities, relationships,
provenance, and neighboring evidence as a compact structured bundle. That
bundle, together with any genuinely unresolved ambiguity, is passed to an
LLM for synthesis. The model's output is validated against the expected
structure. Human review is triggered only when confidence is low, the case
is novel, the risk is high, or the downstream consequence is significant.
Validated outputs update the graph with full provenance.

The result is a system in which the model is invoked selectively, receives
precisely scoped inputs, and produces outputs that are auditable against
structured evidence. Each stage handles what it is best suited for
\parencite{ma2025}.

\subsection{Caching and Reuse}

One of the underappreciated benefits of structured knowledge infrastructure
is the opportunity for caching and reuse. Stable concepts, validated
definitions, known entity relationships, and confirmed schema mappings do
not need to be re-inferred on every query. They can be stored, retrieved
deterministically, and passed to the model as established context rather
than as text to be interpreted.

This matters at scale. An organization that runs thousands of literature
queries per week, each of which reconstructs the same domain terminology
from scratch, is paying a compounding overhead that structured retrieval
eliminates. Token reduction methods that preserve essential meaning while
reducing context size are an active research direction \parencite{trim2024},
but the most reliable form of token reduction is not compression after the
fact; it is not sending information the system already knows.

\subsection{The Human in the Right Place}

Human-in-the-loop is a term that has been overloaded to the point of losing
precision. In many current deployments, it means a human reviewing every
output before it is used, which is operationally expensive and does not
scale. In the architecture described here, it means something more specific:
human expertise applied at the points where it has the highest leverage.

Humans should define the ontology: which concepts are in scope, where
boundaries lie, which synonyms are acceptable, which exclusions matter.
Humans should validate the schema and resolve ambiguous cases that the
system cannot handle deterministically. Humans should review high-impact
outputs, i.e., those with significant downstream consequences, novel cases
outside the system's validated range, and outputs flagged by the validation
layer as uncertain.

Humans do not need to manually inspect every trivial mapping, formatting
correction, or deterministic validation result. The system should triage
review based on confidence, novelty, risk, and downstream consequence
\parencite{garciafernandez2025,tsaneva2025,manzoor2022}. Recent system
designs that place an LLM at the centre of a knowledge-graph
question-answering loop with structured human review at decision points have
made the architectural case explicit: routing model output through
graph-grounded checks and explicit hand-off to domain experts at the points
where evidence is ambiguous or stakes are high is operationally more
reliable than either fully automated or fully manual workflows
\parencite{pusch2026}.

This model places human expertise upstream, where it shapes the system,
rather than downstream, where it repairs its outputs. It is more efficient,
more scalable, and more aligned with how expert judgment actually adds value.

%% ── Section 6 ────────────────────────────────────────────────────────────

\section{Section 6 --- The Organizational Imperative: People, Not Just Architecture}

\subsection{The Talent Gap That Nobody Is Planning For}

Technical architecture documents describe systems. They do not describe
the people needed to build and sustain them. That gap matters, because the
skills required for knowledge infrastructure work are not the same as the
skills required for AI development, and confusing the two leads to
organizational blind spots that are expensive to correct later.

An AI developer or data scientist working with LLMs needs skills in model
selection, prompt engineering, fine-tuning, evaluation, and inference
optimization. These are valuable and increasingly common competencies.

An ontologist or knowledge engineer working on life science data needs
something different. They need to understand the domain: the biology, the
chemistry, the assay design, the therapeutic context, the regulatory
environment. They need to know why ``locomotor activity'' and ``distance
travelled'' are not the same thing, and why that distinction matters when
integrating data across studies. They need to understand metadata standards,
controlled vocabularies, RDF, OWL, SPARQL, and validation frameworks. And
they need to be able to work with domain experts to define concepts,
negotiate boundaries, and maintain a living knowledge system over time.

These are not adjacent skill sets. They are different professions, and they
require different hiring, development, and organizational positioning
strategies.

The gap, however, is narrowing in important ways and should not be framed
as a fixed binary. Hybrid profiles are increasingly visible: computational
biologists who have moved into ontology engineering, bioinformatics-trained
physicians who can operate FHIR or OMOP pipelines, clinical data scientists
with formal ontology training, and AI engineers who have specialized in
biomedical knowledge representation. Graduate programs in biomedical
informatics, computational biology, and digital health are beginning to
produce candidates with substantial cross-coverage of both stacks. The
organizational challenge is therefore less about finding a unicorn than
about (1) recognising which hybrid roles exist and where they sit in the
labour market, (2) positioning them where their leverage is highest---upstream
of system design, not downstream of output correction---and (3) constructing
teams that pair domain-deep knowledge engineers with AI-deep engineers and
give each enough authority to disagree productively with the other.

A note on organizational scale is also warranted. Not every life science
organization can absorb the same investment profile on the same timeline.
Mid-size and small biopharma, biotechs, contract research organizations, and
academic translational groups face genuine capacity constraints: they may not
be able to staff a dedicated ontologist and a dedicated knowledge engineer
and a dedicated AI architect simultaneously, and the schematic USD 30,000 +
USD 500/month figure used in Section 7 is illustrative for a mid-sized
organization rather than universally applicable. For smaller organizations,
three pragmatic paths exist: (a) participate in community efforts such as
the Pistoia Alliance's Pharma General Ontology so that controlled
vocabularies are not reconstructed in-house, (b) adopt established standards
(OMOP CDM for observational analyses, HL7 FHIR for clinical exchange,
established ontologies such as Mondo, ChEBI, NCBI Taxon) rather than
building bespoke ones, and (c) start narrow---instrument one high-value
workflow with structured infrastructure before generalizing. The argument of
this paper is not that every organization must build the same infrastructure
at the same scale; it is that no organization can afford to ignore the
question, and that the cost of ignoring it has shifted.

\subsection{The Market Is Already Signaling This}

The pharmaceutical industry is beginning to recognize this distinction in
practice. Novo Nordisk has publicly described ontology-based data management
as part of its digital transformation in research and early development,
framing structured knowledge infrastructure as operational necessity rather
than academic exercise \parencite{tan2025}. The Pistoia Alliance's Pharma
General Ontology initiative has brought together pharmaceutical stakeholders
to define agreed core entities and recommend controlled terminologies for
data exchange across the industry, explicitly framing this as a community
effort requiring domain expertise, not a problem that model generation can
solve unilaterally \parencite{pistoiaalliance2025}.

At the individual organizational level, large pharmaceutical companies are
already hiring for these hybrid roles. Job descriptions for knowledge graph
engineers in pharma R\&D now routinely specify large-scale knowledge graph
construction, ontology development, pharmaceutical or healthcare domain
integration, and proficiency with RDF, OWL, SPARQL, and graph databases
alongside AI and data science skills \parencite{jj2026}. These roles exist
because the work cannot be done without domain-grounded semantic expertise,
and that expertise cannot be substituted by a larger model.

\subsection{The Dependency Risk}

Organizations that do not develop this capability internally face a specific
and underappreciated risk. It is not only the risk of paying more for AI
inference when prices normalize. It is the risk of having outsourced the
interpretation of their own scientific data to a vendor whose incentives,
architecture choices, and pricing decisions they do not control.

An organization's semantic layer---the concepts, relationships, validation
rules, and metadata standards that govern how its data is understood and
used---is a strategic asset. If that layer is owned by an AI vendor, or
reconstructed on demand by a model the organization does not control, the
organization has not built a data capability. It has rented one.

Vendor lock-in is a known risk in technology procurement, involving
business, technical, and legal dimensions that extend well beyond switching
costs \parencite{oparamartins2016}. In the context of AI, the lock-in is
deeper, because it operates at the level of meaning, not only at the level
of tooling. An organization whose workflows, prompts, semantic assumptions,
agents, and evaluation layers are all tightly coupled to a single model
provider cannot easily disentangle them. The European Parliament's work on
digital sovereignty has framed this concern at a policy level
\parencite{madiega2020}, and European industrial actors are increasingly
treating dependency on non-European technology infrastructure as a security,
reliability, and strategic risk \parencite{chee2025}.

The pragmatic response is not to demand absolute independence from AI
vendors. The Capgemini CEO has argued publicly that full European tech
autonomy is neither achievable nor desirable, given the depth of global
technology interdependencies \parencite{marchandon2026}---though it is worth
noting that Capgemini, as a global consultancy with significant cross-border
AI partnerships, has a direct commercial interest in that position, and it
should be read accordingly.

The fuller picture is more reciprocal. Europe is not only a technology
consumer. It is home to ASML, whose extreme ultraviolet lithography machines
are an irreplaceable component of the global semiconductor supply chain---including
for the chips that power AI infrastructure worldwide. It has produced
Mistral, one of the few non-US frontier model developers operating at scale.
And through GDPR, the AI Act, and its data governance frameworks, it has
become the world's de facto standard-setter in AI regulation, shaping how
organizations on every continent handle data and accountability. The global
AI ecosystem depends on European components, regulation, and scientific
output as much as Europe depends on it.

This mutual dependency does not make the sovereignty argument irrelevant.
It makes it more precise. The question is not whether European life science
organizations should sever ties with US or Asian model providers. It is
whether they retain enough internal capability---in semantic infrastructure,
in knowledge engineering, in domain-grounded data governance---to remain
strategic actors rather than passive consumers. The goal is not autarky. It
is agency. Organizations that do not own their semantic layer will rent not
only compute, but interpretation.

\subsection{Four Questions for Leadership}

For a senior leader in a life science or pharmaceutical organization, the
organizational question reduces to four practical assessments.

First: do you have people who combine domain expertise with knowledge
engineering skills---who understand the biology and the ontology, the
chemistry and the metadata standard, the therapeutic area and the validation
rule? If not, you have a dependency that will become visible when the AI
bill arrives or when a critical workflow fails.

Second: are those people positioned to shape your systems rather than repair
their outputs? The test is simple: did they have a seat at the table when
your current AI workflows were designed, or were they brought in afterwards
to review what the system produced? Upstream expertise and downstream review
are not the same investment.

Third: do you know what your semantic layer costs to replace? Most
organizations that have not built one do not realize how much interpretive
work they are currently outsourcing to models, and therefore cannot price
the dependency. A useful exercise is to ask: if our primary model provider
doubled its prices tomorrow, which workflows would break, which would
degrade, and which would continue unaffected because they rest on structured
knowledge we own? The answers reveal the true shape of the exposure.

Fourth: who in your organization has authority over the conceptual
definitions that govern your data---the terms, the relationships, the
validation rules, the metadata standards? If the answer is ``nobody in
particular'' or ``it depends on the system,'' that is the gap. Data
foundations require governance, and governance requires ownership. Building
on sand is not only a technology risk. It is a leadership gap.

%% ── Section 7 ────────────────────────────────────────────────────────────

\section{Section 7 --- How to Justify Investment in Data Foundations}

\subsection{The Perennial Problem}

The argument for investing in data foundations is not new. It has been
made, in various forms, by data governance practitioners, information
architects, and scientific data managers for decades. It has consistently
struggled to gain traction in budget conversations, for a simple reason: the
foundation is not the application.

When a leadership team reviews an investment proposal, they naturally focus
on the outcome: the search capability, the analytics dashboard, the
literature monitoring tool, the regulatory submission workflow. The
foundation that makes those outcomes possible---the metadata schema, the
ontology, the validation layer, the knowledge graph---is invisible when it
works and blamed for other things when it does not. The case for it has
always been indirect.

What has changed is the cost of not having it. And that cost is now becoming
legible.

\subsection{A New Framing: The Cost of Absence}

The traditional ROI framing for data foundations---``invest now and benefit
later''---has limited persuasive power because the benefits are diffuse and
the timeline is long. A more effective framing in the current environment is
the cost of absence: what does it cost, in AI inference, in human
correction, in re-curation, in missed reuse, to operate without structured
knowledge infrastructure?

The answer, at scale, is substantial. An organization running AI-powered
literature review, entity extraction, data integration, and regulatory
reporting across a large corpus, without metadata standards or ontologies to
reduce context and pre-filter retrieval, is paying a compounding overhead on
every query. That overhead is not visible as a line item. It is distributed
across inference costs, human review time, error correction, and the
opportunity cost of results that are slower, less reliable, and harder to
audit than they would be with structure in place.

\subsection{A Schematic Cost Scenario}

To make the cost-of-absence framing concrete, consider a schematic scenario
for a mid-sized life science organization. The scenario is illustrative, not
predictive: the precise figures depend on team size, model selection,
pricing, and the structure of the underlying corpus. Each value below is a
rounded planning number, not a measured outcome. We present it to show how
the categories compose, not as a forecast.

Assume an organization running roughly \textbf{1,000 AI-assisted
scientific-context queries per week}---literature triage, protocol
comparison, regulatory summarization, entity extraction, and similar tasks.
Assume each unstructured query draws on the equivalent of two documents
pulled into the context window, where a structured-retrieval variant of the
same query would draw on a compact subgraph plus targeted snippets. Under
current frontier pricing, and across a representative mix of task types, the
structured approach can reduce per-query inference cost by \textbf{roughly
one order of magnitude} for the inference component, with additional savings
on human review time because outputs become more auditable and validation can
be partially automated.

Set against this, assume an infrastructure investment of approximately
\textbf{USD 30,000} to stand up a usable initial knowledge graph and
metadata schema (covering tooling, data modelling, and initial population
from existing curated sources), plus approximately \textbf{USD 500 per
month} in ongoing operating cost (hosting, validation, light curation;
excludes the fully loaded cost of dedicated knowledge engineering staff,
which is a separate hiring decision discussed in Section 6). Under those
assumptions, the inference-cost differential alone produces a \textbf{break-even
point of approximately 14--15 months}, with several caveats.

The caveats matter. The order-of-magnitude inference reduction depends on
task mix; tasks that are intrinsically open-ended or require frontier
reasoning will not benefit as much as structured retrieval and validation
tasks. The USD 30,000 figure assumes the organization has existing curated
data sources (controlled vocabularies, study metadata, internal taxonomies)
that can be ingested with moderate engineering effort; from a true cold
start, both the cost and the timeline rise. The USD 500/month operating
estimate excludes the cost of a dedicated knowledge engineer or ontologist;
a serious deployment includes one. And, critically, the break-even
calculation captures only the inference-cost component; it does not capture
the larger and more durable benefits of auditability, reuse across projects,
reduced human correction load, and reduced vendor dependency, which compound
over a longer horizon.

The point of the scenario is not the specific 14--15 month figure. It is
that the cost of \textit{not} building structured knowledge infrastructure,
evaluated against current and announced pricing trajectories, is no longer
obviously cheaper than the cost of building it---and at the scale at which
AI workflows are now being deployed across pharmaceutical R\&D, the
calculation increasingly tilts toward investment. Organizations should
perform this calculation with their own numbers; the categories are stable
even where the magnitudes vary.

Data without appropriate metadata cannot be fully interrogated or integrated
into new projects, creating wasted resources and missed opportunities for
reuse \parencite{moresis2024}. FAIR data principles make explicit that the
goal of machine-actionability is not an aspiration; it is a prerequisite for
automated systems to operate reliably at scale \parencite{wilkinson2016}.
Community-defined metadata standards that encode domain requirements have
been shown to support verifiable quality, interoperability, and complex
reporting needs \parencite{batista2022,musen2022}.

\subsection{The Portfolio Argument}

A second reframing that resonates with leadership is the portfolio argument
\parencite{vereckey2025}. Data foundations are rarely justified by the first
application they enable. They are justified by the portfolio of applications
that become possible over time, and by the reduction in per-application cost
that comes from having the foundation already in place.

An ontology built to support one preclinical data integration project does
not need to be rebuilt for the next one. A metadata schema validated for one
regulatory submission workflow can be extended for the next therapeutic area.
A knowledge graph populated with one year's experimental data becomes more
valuable with each year that follows, because the number of queryable
relationships grows, the number of comparison points increases, and the cost
of answering new questions falls.

Leading organizations focus their AI efforts more narrowly and expect higher
ROI than peers precisely because they invest in the infrastructure that makes
multiple applications efficient, rather than building each application from
scratch \parencite{apotheker2025}. The ROI of data foundations is often
indirect because the foundation is not the application. It shows up as
reduced duplication, faster reuse, better retrieval, lower validation burden,
improved auditability, lower inference waste, and higher probability that AI
systems can be deployed safely at scale.

The converse is equally relevant: only about 25\% of AI initiatives deliver
their expected ROI, and only 16\% scale enterprise-wide \parencite{ibm2025}.
Weak data foundations and weak governance are consistently identified as
contributing factors. The investment in structured knowledge infrastructure
is, among other things, an investment in the probability that AI initiatives
succeed.

\subsection{Connecting to Responsible AI Scaling}

There is a final dimension to the business case that is increasingly
relevant for organizations with public sustainability and governance
commitments. Data centre electricity demand is growing sharply, driven in
significant part by AI workloads \parencite{noffsinger2025}. Reducing
unnecessary inference through better knowledge infrastructure is not a niche
technical optimization. It is part of responsible AI scaling, a way of
ensuring that the organization's AI investment generates value proportional
to its environmental and financial cost.

For pharmaceutical and life science organizations, this argument has
particular force. The regulatory environment around AI in drug development
is evolving rapidly, and auditability, provenance, and reproducibility are
not optional features. They are requirements. The FDA's evolving guidance on
AI/ML-enabled medical devices and on the use of AI in regulatory
decision-making for drugs and biological products emphasises lifecycle
traceability, documented data provenance, and pre-specified evaluation
\parencite{fda2025,fda2024}. The European Medicines Agency's reflection
paper on the use of AI in the medicinal product lifecycle is closely aligned:
it calls for documented data quality, model transparency, and reproducible
analytical pipelines from preclinical through post-authorisation phases
\parencite{ema2024}. The EU AI Act applies high-risk obligations---including
data governance, technical documentation, traceability, and human
oversight---to AI systems used in regulated medical contexts, and those
obligations bind organizations regardless of where the underlying model was
trained \parencite{euaiact2024}. Structured knowledge infrastructure is not
in tension with these requirements. It is how they are met.

%% ── Section 8 ────────────────────────────────────────────────────────────

\section{Section 8 --- Why This Matters in Life Sciences}

\subsection{The Particular Stakes of Getting It Wrong}

In most domains, a hallucinated output or a poorly integrated dataset is an
inconvenience. In life sciences, it can be a safety issue, a regulatory
failure, or years of wasted research investment. The stakes of getting the
knowledge infrastructure right are not abstract.

Preclinical research produces large volumes of experimental data across
species, models, protocols, assay types, and therapeutic areas. That data
is only useful at scale if it can be compared, integrated, and reused across
studies. Comparison requires shared terminology. Integration requires
compatible metadata. Reuse requires provenance. Without these, each dataset
is an island, legible to the team that produced it and opaque to everyone
else, including the AI systems that are increasingly being asked to extract
insight from it.

\subsection{Metadata and Reproducibility}

The connection between metadata standards and scientific reproducibility is
well established. ARRIVE 2.0 defines the minimum reporting information for
animal research, connecting standardized metadata to reproducibility and
reporting quality, not only to AI \parencite{dusert2020}. A minimal
metadata set designed specifically to support repurposing of nonclinical in
vivo data has been proposed and validated in the literature, explicitly
aligned with ARRIVE 2.0 and FAIR compliance \parencite{moresis2024}. On the
clinical side, the OMOP Common Data Model and HL7 FHIR play a parallel
structural role: OMOP supports reproducible observational and real-world
evidence analyses across institutions \parencite{hripcsak2015}, while FHIR
underpins interoperable exchange of clinical and trial data across systems
and regulators \parencite{bender2013}. Both exemplify how machine-actionable
structure earns its return in regulated biomedical contexts independent of
any particular AI capability. The argument is not that metadata is needed
because AI requires it. It is that metadata is needed because science
requires it, and AI makes the absence of it more expensive.

\subsection{Knowledge Graphs in Biomedical Research}

The life sciences have been among the most active domains for knowledge
graph development, for reasons that are directly related to the complexity
of the domain's conceptual landscape. Biomedical concepts are numerous,
interrelated, often ambiguous across contexts, and organized by multiple
competing ontologies that must be reconciled for data integration to work.

Biomedical knowledge graph ecosystems that integrate ontologies and
heterogeneous life science data are being developed explicitly for
AI-powered research, combining experimental results, literature, pathway
data, and clinical information in a structured form that supports both
querying and model-based reasoning \parencite{callahan2024}. PubMed
knowledge graph systems have demonstrated the strategic value of connecting
papers, patents, clinical trials, biomedical entities, citations, and author
networks at scale \parencite{xu2025}. Graph data models are increasingly used
to structure biomedical and clinical information in ways that enable new
forms of analysis over heterogeneous data that relational databases cannot
easily support \parencite{hansel2023}.

PubMed metadata converted to a knowledge graph format has been shown to
support semantic biomedical retrieval using entities, MeSH terms, citations,
grants, and author metadata, demonstrating the practical value of the
metadata-to-graph-to-retrieval pipeline in a domain directly relevant to
pharmaceutical research \parencite{ebeid2022}.

\subsection{Ontologies as Operational Infrastructure}

For a long time, ontologies in pharma were associated primarily with
academic bioinformatics and standards bodies. That association is changing.
The Novo Nordisk case is instructive: a major pharmaceutical company has
publicly framed ontology-based data management as part of its core digital
transformation in research and early development, not as a future aspiration
but as current operational practice \parencite{tan2025}.

The Pistoia Alliance's Pharma General Ontology initiative represents a
broader industry recognition that interoperability across FAIR datasets
requires shared, community-maintained controlled vocabularies, and that
building those vocabularies requires the kind of domain-grounded expertise
described in Section 6 \parencite{pistoiaalliance2025}.

FAIR principles recommend controlled vocabularies and ontologies to define
data and metadata concepts, and this is not a recommendation made in the
context of AI adoption. It predates the current generation of AI tools and
reflects a longstanding understanding in biomedical data management that
semantic precision is a prerequisite for scientific reuse
\parencite{bernabe2023}.

\subsection{The Specific Case for Preclinical Data}

The preclinical domain illustrates the argument with particular clarity.
Consider the challenge of constructing a virtual control group---drawing on
historical control data to reduce the number of animals used in new studies.
This requires identifying historical experiments that are sufficiently
comparable on species, strain, age, sex, housing conditions, route of
administration, assay protocol, and endpoint definition. Without standardized
metadata encoding those variables consistently, the comparison cannot be made
reliably. The AI system, however capable, cannot infer comparability from
unstructured text with sufficient precision for a regulatory context.

With standardized metadata, ontology-mapped endpoints, and a knowledge graph
that preserves provenance and experimental conditions, the comparison becomes
a structured query. The model's role is reduced to synthesis and
interpretation of a pre-filtered, structurally coherent evidence set. The
result is faster, cheaper, more reproducible, and more auditable---all of
which matters in an environment where regulators are increasingly
scrutinizing the provenance and reliability of AI-assisted analysis.

%% ── Section 9 ────────────────────────────────────────────────────────────

\section{Section 9 --- Risks, Limitations, and What This Approach Does Not Solve}

\subsection{Intellectual Honesty About the Constraints}

Any argument for a particular architectural approach is incomplete without
an honest account of its limitations. The case for graphs, ontologies, and
metadata standards is strong, but it is not a case for a frictionless
solution. These components have real costs and genuine failure modes that
organizations need to understand before investing.

\subsection{Graphs Are Not Free}

Knowledge graphs require design, curation, tooling, and maintenance.
Building a graph that accurately represents a complex domain is a
significant engineering and intellectual undertaking. Populating it requires
either manual curation, which is expensive and slow, or automated
extraction, which introduces its own error rates. Maintaining it as the
domain evolves, as new assay types emerge, as regulatory standards change,
as scientific understanding advances, requires ongoing investment.

Ontology-guided extraction with in-context learning can support
semi-automated knowledge graph construction for domain-specific data with
limited labelled examples \parencite{vancauter2024}, and LLMs can assist in
drafting ontology structures from domain descriptions \parencite{lippolis2025}.
But these are aids, not substitutes, for the domain expertise and
engineering discipline that knowledge graph construction requires. An
organization that expects to deploy a production-quality biomedical knowledge
graph without sustained investment in people, tooling, and process will be
disappointed.

\subsection{Ontologies Can Become Too Rigid}

An ontology that is designed without deep involvement from domain users
risks encoding the wrong distinctions, or encoding them too rigidly for the
actual variety of the domain. Ontologies that are too fine-grained become
difficult to maintain and fail to generalize; those that are too
coarse-grained fail to capture the distinctions that matter for data
integration.

LLM support for ontology engineering is an active research area, but it
remains fragmented and benchmark-challenged \parencite{garijo2025}. LLMs
can support ontology extension and drafting, but ontology engineering still
requires the kind of human-LLM collaboration in which domain experts
validate conceptual boundaries and relations that automated generation
cannot reliably produce \parencite{garciafernandez2025}. The risk of
delegating ontology design to a model without expert validation is an
ontology that looks formally correct but encodes semantic errors that
propagate through every downstream system.

\subsection{GraphRAG Can Still Hallucinate}

Grounding a model's output in a knowledge graph reduces hallucination risk
but does not eliminate it. If the retrieved subgraph is wrong, incomplete,
or poorly linearized into the model's context, the model can still produce
inaccurate outputs. KG-LLM integration remains an active research area with
unresolved challenges around datasets, benchmarks, knowledge integration,
and hallucination evaluation \parencite{lavrinovics2024}. The claim in this
paper is that graphs give you better tools to detect and audit hallucination;
it is not that they prevent it.

\subsection{Bad Metadata Creates False Precision}

A metadata standard that is poorly designed, inconsistently applied, or
maintained without quality control can be worse than no standard at all,
because it creates the illusion of structured, comparable data where no real
comparability exists. Token reduction should be measured empirically, not
assumed \parencite{trim2024}. The benefits of structured knowledge
infrastructure are contingent on the quality of that structure. Quality
requires investment, governance, and ongoing validation.

\subsection{Human Review Can Become a Bottleneck}

The human-in-the-loop model described in Section 5 depends on human review
being applied selectively, at the points of highest leverage. If review is
applied too broadly, the bottleneck moves from the model to the reviewer,
and the efficiency gains of the structured pipeline are consumed by the cost
of manual oversight. The system must be designed to triage review
intelligently, surfacing only the cases where human judgment is genuinely
needed \parencite{tsaneva2025}.

\subsection{Why the Argument Holds Under Optimistic Model Improvement}

A reasonable objection to the case we have made is that continued model
improvement may dissolve some of the costs we describe: a sufficiently
capable future model might infer missing context, reconstruct provenance,
harmonize terminology across studies, and self-validate its outputs to the
point where much of the infrastructure we advocate becomes redundant. This
objection deserves a direct response, not a dismissal.

We argue the investment case holds even under optimistic assumptions about
model trajectory, for three reasons that are structural rather than
capability-bound.

First, \textbf{auditability and versioning are regulatory requirements, not
capability requirements}. In regulated life science contexts, what matters
is not only that an answer is correct but that the chain of evidence and
definitional choices behind it is independently inspectable, versioned, and
reproducible by a third party. A model that infers context accurately on a
single query does not by itself produce an audit-grade record of \textit{how}
it inferred it, against \textit{which} definitional baseline, and
\textit{whether} that baseline matches the one the regulator or the next
investigator will use. Structured knowledge infrastructure exists in part to
externalize this audit surface from the model. Better model reasoning does
not replace that externalization; it operates against it.

Second, \textbf{interoperability requires stable schemas that no individual
model can unilaterally guarantee}. Data integration across studies, sites,
sponsors, and decades depends on shared, durable, community-maintained
representations of meaning. A model that reads two unstructured datasets and
``understands'' both is not the same as two datasets that conform to a
shared ontology and a common metadata standard. Only the latter supports
reliable downstream integration by \textit{other} systems, \textit{other}
models, and \textit{other} organizations operating on different software
stacks and at different points in time. Interoperability is a property of
the data layer, not of any particular model.

Third, \textbf{better models tend to \textit{increase}, not decrease, the
marginal value of structured knowledge infrastructure}. A capable model
paired with a high-quality ontology, knowledge graph, and validation layer
can do strictly more than the same model operating on unstructured text: it
can decompose tasks more precisely, retrieve more focused evidence, attribute
claims to specific sources, and route subtasks deterministically where
appropriate. The complementarity is asymmetric in favour of structure:
capable models amplify the leverage of good infrastructure more than they
substitute for it. Organizations that under-invest in structure will not
catch up to better infrastructure simply by buying access to better models.

These three reasons are not contingent on assumptions about whether the next
model generation will or will not be substantially more capable. They
concern, respectively, what regulators require, what cross-organizational
data exchange requires, and how capability and structure compose at scale.

\subsection{Failure Conditions: When This Argument Would Weaken}

Intellectual honesty also requires being explicit about the conditions under
which the case we have made would weaken. The argument we make is not
unconditional. It would lose force, in roughly increasing order of how
plausible we believe each condition to be on a 3--5 year horizon, if models
could reliably and verifiably do \textit{all five} of the following:

\begin{enumerate}
  \item \textbf{Infer all relevant experimental and clinical context} from
  primary documents, including unstated assumptions, protocol amendments,
  instrument calibration history, and tacit lab conventions, to a standard
  acceptable in a regulatory submission, without prior structuring.

  \item \textbf{Preserve and reconstruct provenance automatically},
  producing for every claim a machine-readable, third-party-verifiable trace
  back to source documents, dataset versions, and experimental conditions,
  without an externalized provenance layer.

  \item \textbf{Align with community standards autonomously}, mapping any
  organization's internal terminology to evolving consensus vocabularies
  (e.g., OMOP CDM, HL7 FHIR, FAIR-aligned ontologies) without explicit
  ontology engineering, and remaining aligned as those standards change.

  \item \textbf{Expose calibrated, machine-actionable uncertainty} at the
  level of individual claims, sufficient for automated downstream systems to
  route low-confidence outputs to human review or alternative evidence
  sources without bespoke validation logic.

  \item \textbf{Satisfy regulatory audit requirements} for AI-assisted
  analysis without any prior structuring of inputs---i.e., demonstrating to
  a regulator that the model's reasoning process and evidence selection are
  reproducible, inspectable, and stable across reruns.
\end{enumerate}

If \textit{all five} of these conditions were met, the marginal value of
organization-owned structured knowledge infrastructure would fall
substantially. We do not believe this is achievable on the 3--5 year horizon
to which this paper's recommendations apply, but the test is empirical, not
rhetorical: the conditions above are operational, observable, and falsifiable.

\subsection{Threshold Analysis: What Would Tell Us We Are Approaching the Failure Conditions}

Beyond stating the failure conditions, it is useful to identify the
empirical signals that would indicate we are approaching them. For each
condition, observable thresholds exist.

For \textbf{context inference} (condition 1), the relevant signal is
independent benchmark evidence that frontier models, evaluated on
representative samples of preclinical or clinical primary documents, can
recover the variables required for a virtual control group reconstruction
(species, strain, age, sex, route of administration, assay protocol,
endpoint definition) with regulatory-grade accuracy and \textit{without}
prior schema-mapping. At present, the evidence base shows degradation on
long contexts, lost-in-the-middle behaviour, and inconsistency across model
families \parencite{liu2024,leng2024}; we are not close to the threshold.

For \textbf{automatic provenance} (condition 2), the threshold signal would
be machine-checkable provenance traces produced by models without external
scaffolding, validated against established provenance standards. Current
systems produce provenance only when prompted or when integrated with
external provenance layers; the gap is structural, not graduating.

For \textbf{community-standard alignment} (condition 3), the signal would
be reproducible, audited mappings from arbitrary organizational terminology
to community-maintained vocabularies (OMOP, FHIR, Mondo, ChEBI, etc.)
produced by models without bespoke ontology engineering. Current evidence
indicates LLMs can \textit{assist} ontology engineering but cannot reliably
replace it \parencite{garciafernandez2025,garijo2025}.

For \textbf{calibrated machine-actionable uncertainty} (condition 4), the
threshold is reliable, well-calibrated confidence scores at the claim level,
validated against held-out ground truth across representative tasks.
Calibration evaluation remains an active research area with substantial gaps.

For \textbf{regulatory audit acceptance} (condition 5), the threshold is
explicit acceptance by major regulators (FDA, EMA, PMDA) of AI-assisted
analyses where the input has not been pre-structured. Current and emerging
regulatory positions (FDA AI/ML SaMD guidance, EMA reflection paper on AI
in the medicinal product lifecycle, EU AI Act high-risk provisions) move in
the opposite direction: they emphasise auditability, traceability, and
pre-specified data structuring as prerequisites for trustworthy AI use in
regulated contexts \parencite{fda2025,ema2024,euaiact2024}.

In short: each failure condition can be operationally monitored. A reader
who suspects this paper's investment case is fragile to model improvement
can answer that suspicion empirically, by tracking the five thresholds
above. We argue the cumulative weight of current evidence does not support
abandoning structured knowledge infrastructure on the 3--5 year horizon,
but we welcome that determination being made on evidence rather than on
faith in either direction.

%% ── Section 10 ───────────────────────────────────────────────────────────

\section{Section 10 --- Conclusion: Constraint as a Competitive Advantage}

\subsection{What We Have Argued}

This paper has made a case in two registers, one strategic and one
technical, that we hope are mutually reinforcing rather than merely parallel.

The strategic argument is this: the early experience of LLMs created a
misleading impression about the cost and complexity of AI-powered knowledge
work. The pricing was subsidized, the dependency was underpriced, and the
structural limitations of unguided model use were invisible at small scale.
That is changing. As AI workflows become embedded in organizational
processes, as pricing shifts toward usage-sensitive models, and as the cost
of inference, energy, and human correction becomes legible, the organizations
that have invested in knowledge infrastructure will be significantly better
positioned than those that have not.

The technical argument is this: LLMs are good at language, not at memory,
meaning, or discipline. Graphs are good at memory. Ontologies are good at
meaning. Metadata standards are good at discipline. Asking one component to
do all four jobs produces a system that does all four jobs expensively and
inconsistently. The right architecture assigns each job to the component
best suited for it, positions the model as an orchestrated component rather
than a universal solution, and places human expertise upstream, where it
shapes the system, rather than downstream, where it repairs its outputs.

\subsection{The Organizational Synthesis}

The dimension of this argument that we believe has received insufficient
attention in the existing literature is the organizational one.
Architectural diagrams and technical papers describe systems. They do not
describe the people needed to build and sustain them, the organizational
structures needed to support those people, or the budget conversations needed
to justify the investment.

The people who build and maintain knowledge infrastructure in life sciences
are domain experts first. Their value comes from understanding the biology,
the chemistry, the regulatory context, and the data well enough to make
precise conceptual decisions. That expertise cannot be substituted by a
larger model or a smarter prompt. It requires hiring, development,
organizational positioning, and sustained investment.

Organizations that do not plan for this talent will find, when the pricing
shift fully arrives, that they have built AI capabilities on a foundation
they do not own. They will be dependent on vendors not only for compute, but
for interpretation. The semantic layer of their data---the part that encodes
what their data means and how it can be used---will be reconstructed on
demand, expensively and imprecisely, by systems they do not control.

\subsection{Three Things to Do in the Next Twelve Months}

For a senior leader in a life science or pharmaceutical organization reading
this paper, the practical implications reduce to four priorities.

First: audit your semantic layer. Do you have ontologies, metadata standards,
and validation rules that govern your key data assets? Are they
machine-actionable, domain-validated, and actively maintained? If not,
identify the highest-priority gaps and begin closing them, starting with the
data that your AI systems rely on most heavily.

Second: assess your talent. Do you have people who combine domain expertise
with knowledge engineering skills? Are they positioned to shape your data
systems, or are they reviewing outputs they had no role in designing? If the
role does not exist in your organization, consider what it would take to
build it, whether through hiring, development, or partnership with academic
groups or industry consortia such as the Pistoia Alliance.

Third: reframe your AI investment thesis. The question is not which AI tools
to buy. It is what knowledge infrastructure those tools require to deliver
reliable, auditable, reusable results, and whether that infrastructure is
being built in proportion to the AI investment. The return on frontier model
access is limited by the quality of the inputs it receives.

Fourth: assign governance ownership. Semantic infrastructure without
governance degrades. Decide who in your organization has authority over the
conceptual definitions, validation rules, and metadata standards that govern
your key data assets. If that authority is diffuse or absent, the
infrastructure you build will not hold. This is not a technical decision. It
is a leadership one.

\subsection{The Final Argument}

Constraint is not a limitation. It is how organizations stay in control of
cost, quality, and knowledge.

An unconstrained LLM workflow feels flexible. Much of that flexibility is
paid for through ambiguity, retries, hallucination risk, and human
correction. A constrained workflow---one in which metadata standards,
ontologies, graphs, validators, and human expertise each operate at the
point of highest leverage---trades that apparent flexibility for something
more valuable: reliability, auditability, reusability, and the ability to
know, with confidence, what your AI systems actually know.

Organizations that invest in that infrastructure now are not choosing caution
over ambition. They are building the foundation on which AI-powered knowledge
work, at scale, in regulated environments, with scientific and legal
accountability, actually rests.

The organizations that do not own their semantic layer will rent not only
compute, but interpretation. The organizations that do will have built
something that no vendor can take away.

%% ── References ───────────────────────────────────────────────────────────

\printbibliography

\end{document}
